\section{Konstanten und Einheiten}
Hier sind Konstanten und ihren verschiedenen Einheiten, Einheiten von variablen und wichtige Umrechnungen aufgelistet.

\subsection{Konstanten}

% --- KONSTANTEN: WIEDER AUFGETEILT ABER EXAKT BÜNDIG ---
% Spalte 1: X (flexibel, mit \hspace{0pt} für Wortumbruch)
% Spalte 2: p{1.2cm} (Feste Breite, erzwingt gleiche Ausrichtung aller Tabellen)
% Spalte 3: X (flexibel für die Werte)

\begin{center}
	\renewcommand{\arraystretch}{1.3}
	\begin{tabularx}{\linewidth}{|>{\raggedright\arraybackslash\hspace{0pt}}X|>{\centering\arraybackslash}p{1.2cm}|>{\raggedright\arraybackslash}X|}
		\hline
		\multicolumn{3}{|c|}{\textbf{Fundamentale Naturkonstanten (SI-Definitionen)}} \\ \hline
		Licht\-geschwindig\-keit im Vakuum & $c$ & $299\,792\,458\,\mathrm{m/s}$ (ca. $3 \cdot 10^8\,\mathrm{m/s}$) \\ \hline
		Elektrische Elementarladung & $e$ & $1{,}602176634 \cdot 10^{-19}\,\mathrm{C}$ \\ \hline
		Planck'sches Wirkungsquantum & $h$ & $6{,}62607015 \cdot 10^{-34}\,\mathrm{J \cdot s}$ \\ \hline
		Boltzmann-Konstante & $k_{\mathrm{B}}$ & $1{,}380649 \cdot 10^{-23}\,\mathrm{J/K}$ \\ \hline
		Avogadro-Konstante & $N_{\mathrm{A}}$ & $6{,}02214076 \cdot 10^{23}\,\mathrm{mol^{-1}}$ \\ \hline
		Frequenz des Hyperfein\-struktur\-übergangs ($^{133}\mathrm{Cs}$) & $\Delta\nu$ & $9\,192\,631\,770\,\mathrm{s^{-1}}$ \\ \hline
		Photometrisches Strahlungsäquivalent & $K_{\mathrm{cd}}$ & $683\,\mathrm{lm/W}$ \\ \hline
	\end{tabularx}
\end{center}

\vspace{-4mm}

\begin{center}
	\renewcommand{\arraystretch}{1.3}
	\begin{tabularx}{\linewidth}{|>{\raggedright\arraybackslash\hspace{0pt}}X|>{\centering\arraybackslash}p{1.2cm}|>{\raggedright\arraybackslash}X|}
		\hline
		\multicolumn{3}{|c|}{\textbf{Elektromagnetismus \& Gravitation}} \\ \hline
		Elektrische Feldkonstante (Dielektrizitäts\-konstante) & $\epsilon_0$ & $8{,}854 \cdot 10^{-12}\,\mathrm{C^2/(N \cdot m^2)}$ \\ \hline
		Proportionalitäts\-faktor (Coulomb-Gesetz) & $1 / (4\pi\epsilon_0)$ & $8{,}99 \cdot 10^9\,\mathrm{N \cdot m^2/C^2}$ \\ \hline
		Magnetische Feldkonstante (Permeabilität des Vakuums) & $\mu_0$ & $4\pi \cdot 10^{-7}\,\mathrm{T \cdot m/A}$ (oder $\mathrm{N/A^2}$) \\ \hline
		Universelle Gravitations\-konstante & $G$ (bzw. $\Gamma$) & $6{,}67 \cdot 10^{-11}\,\mathrm{N \cdot m^2/kg^2}$ \\ \hline
	\end{tabularx}
\end{center}

\vspace{-4mm}

\begin{center}
	\renewcommand{\arraystretch}{1.3}
	\begin{tabularx}{\linewidth}{|>{\raggedright\arraybackslash\hspace{0pt}}X|>{\centering\arraybackslash}p{1.2cm}|>{\raggedright\arraybackslash}X|}
		\hline
		\multicolumn{3}{|c|}{\textbf{Thermodynamik \& Atome/Teilchen}} \\ \hline
		Universelle Gaskonstante & $R$ & $8{,}314\,\mathrm{J/(mol \cdot K)}$ \\ \hline
		Masse eines Elektrons & $m_{\mathrm{e}}$ & $9{,}11 \cdot 10^{-31}\,\mathrm{kg}$ \\ \hline
		Masse eines Protons & $m_{\mathrm{p}}$ & $1{,}67 \cdot 10^{-27}\,\mathrm{kg}$ \\ \hline
		Bohr'sches Magneton & $\mu_{\mathrm{Bohr}}$ & $9{,}72 \cdot 10^{-24}\,\mathrm{A \cdot m^2}$ (bzw. $5{,}79 \cdot 10^{-5}\,\mathrm{eV/T}$) \\ \hline
	\end{tabularx}
\end{center}


\subsection{Einheiten}

% --- EINHEITEN: AUFGETEILT UND BÜNDIG ---
% Spalte 1: X (flexibel)
% Spalte 2 & 3: p{2.2cm} (Fest, damit Einheit und Symbol immer exakt gleich breit sind)

\begin{center}
	\renewcommand{\arraystretch}{1.3}
	\begin{tabularx}{\linewidth}{|>{\raggedright\arraybackslash\hspace{0pt}}X|>{\raggedright\arraybackslash}p{2.2cm}|>{\raggedright\arraybackslash}p{2.2cm}|}
		\hline
		\textbf{Physikalische Größe} & \textbf{Einheit} & \textbf{Symbol} \\ \hline
		\multicolumn{3}{|c|}{\textbf{SI-Basiseinheiten}} \\ \hline
		Länge & Meter & $\mathrm{m}$ \\ \hline
		Masse & Kilogramm & $\mathrm{kg}$ \\ \hline
		Zeit & Sekunde & $\mathrm{s}$ \\ \hline
		Elektrischer Strom & Ampere & $\mathrm{A}$ \\ \hline
		Thermodynamische Temperatur & Kelvin & $\mathrm{K}$ \\ \hline
		Stoffmenge & Mol & $\mathrm{mol}$ \\ \hline
		Lichtstärke & Candela & $\mathrm{cd}$ \\ \hline
	\end{tabularx}
\end{center}

\vspace{-4mm}

\begin{center}
	\renewcommand{\arraystretch}{1.3}
	\begin{tabularx}{\linewidth}{|>{\raggedright\arraybackslash\hspace{0pt}}X|>{\raggedright\arraybackslash}p{2.2cm}|>{\raggedright\arraybackslash}p{2.2cm}|}
		\hline
		\multicolumn{3}{|c|}{\textbf{Mechanik, Thermodynamik \& Optik}} \\ \hline
		Kraft & Newton & $\mathrm{N} = \mathrm{kg \cdot m/s^2}$ \\ \hline
		Druck & Bar / $\mathrm{N/m^2}$ & $\mathrm{bar}$ / $\mathrm{N/m^2}$ \\ \hline
		Energie / Arbeit & Joule & $\mathrm{J} = \mathrm{N \cdot m} = \mathrm{kg \cdot m^2/s^2}$ \\ \hline
		Leistung & Watt & $\mathrm{W} = \mathrm{J/s}$ \\ \hline
		Frequenz & Hertz & $\mathrm{Hz} = \mathrm{s^{-1}}$ \\ \hline
		Ebener Winkel & Radiant & $\mathrm{rad}$ \\ \hline
		Raumwinkel & Steradiant & $\mathrm{sr}$ \\ \hline
		Volumen & Liter & $\mathrm{l}$ oder $\mathrm{L} = 10^{-3}\,\mathrm{m^3}$ \\ \hline
		Länge (Astronomie) & Lichtjahr & $\mathrm{Lj}$ \\ \hline
		Länge (Astronomie) & Astr. Einheit & $\mathrm{AE}$ \\ \hline
		Zeit (Alltag) & Min., Std., Tag, Jahr & $\mathrm{min}$, $\mathrm{h}$, $\mathrm{d}$, $\mathrm{a}$ \\ \hline
		Temperatur (Alltag) & Grad Celsius & $^{\circ}\mathrm{C}$ \\ \hline
	\end{tabularx}
\end{center}

\vspace{-4mm}

\begin{center}
	\renewcommand{\arraystretch}{1.3}
	\begin{tabularx}{\linewidth}{|>{\raggedright\arraybackslash\hspace{0pt}}X|>{\raggedright\arraybackslash}p{2.2cm}|>{\raggedright\arraybackslash}p{2.2cm}|}
		\hline
		\multicolumn{3}{|c|}{\textbf{Elektromagnetismus}} \\ \hline
		Elektrische Ladung & Coulomb & $\mathrm{C} = \mathrm{A \cdot s}$ \\ \hline
		Spannung & Volt & $\mathrm{V} = \mathrm{J/C}$ \\ \hline
		Elektrische Kapazität & Farad & $\mathrm{F} = \mathrm{C/V}$ \\ \hline
		Elektrischer Widerstand & Ohm & $\Omega = \mathrm{V/A}$ \\ \hline
		Elektrische Feldstärke & Volt pro Meter & $\mathrm{V/m}$ oder $\mathrm{N/C}$ \\ \hline
		Magnetfeldstärke & Tesla & $\mathrm{T} = \mathrm{V \cdot s/m^2}$ \\ \hline
		Magnetfeldstärke (CGS) & Gauß & $\mathrm{G} = 10^{-4}\,\mathrm{T}$ \\ \hline
		Magnetischer Fluss & Weber & $\mathrm{Wb} = \mathrm{V \cdot s}$ \\ \hline
		Induktivität & Henry & $\mathrm{H}$ \\ \hline
		Magnetisches Moment & Ampere-$\mathrm{m^2}$ & $\mathrm{A \cdot m^2}$ \\ \hline
		Energie (Atomphysik) & Elektronenvolt & $\mathrm{eV}$ \\ \hline
	\end{tabularx}
\end{center}